% ==============================================================
%                        CHAPTER 4
% ==============================================================

\chapter{Implementation and Realization}

\section{Introduction}

This chapter is dedicated to the realization of the Stag.io platform. After presenting the conceptual and design models in the previous chapters, we now transition to the concrete technical implementation of the system. We present the proposed system architecture, detailing the hardware and software resources mobilized for development. We then discuss the implementation of the core algorithmic components --- the SmartMatch recommendation engine, the real-time notification system, and the security mechanisms --- supported by actual source code excerpts drawn from the project repository. We also describe the testing strategy used to validate key operations. Finally, we showcase the end functional result through a selection of user interfaces illustrating the software's manipulation.

\section{Proposed System Architecture}

To ensure high availability, scalability, and maintainability, Stag.io was designed using a modern multi-tiered architecture that separates the client interface from the backend processing and data persistence layers. This separation of concerns allows each layer to evolve, scale, and be maintained independently.

\subsection{Component Architecture}

The component diagram illustrates the high-level architecture of Stag.io, detailing the various subsystems, their boundaries, and how they interact to provide the platform's functionalities. It distinctly separates the React frontend interface, the Express backend service layer, and external cloud dependencies like Google Gemini and Cloudflare.

\begin{figure}[H]
\centering
\includegraphics[width=\textwidth,height=0.88\textheight,keepaspectratio]{6-component-diagram.png}
\caption{Component diagram detailing the functional architecture}
\label{fig:component-diagram}
\end{figure}

\subsection{Deployment Architecture}

The deployment diagram provides a physical view of the system, showing how the software components are distributed across the cloud infrastructure. It maps the Next.js client environment to edge servers, the Node.js backend to application servers, and the database to a managed serverless persistence layer.

\begin{figure}[H]
\centering
\includegraphics[width=\textwidth,height=0.88\textheight,keepaspectratio]{7-deployment-diagram.png}
\caption{Deployment diagram showing node distribution in Stag.io}
\label{fig:deployment-diagram}
\end{figure}

\section{Hardware and Software Resources}

The realization of this project required a carefully selected set of hardware devices and software frameworks.

\subsection{Hardware Configuration}

Table~\ref{tab:hardware} summarizes the hardware environment used during development and deployment.

\begin{table}[H]
\centering
\caption{Hardware resources used for Stag.io}
\label{tab:hardware}
\begin{tabular}{|p{3.5cm}|p{10.5cm}|}
\hline
\textbf{Resource} & \textbf{Specification} \\
\hline
Processor & Intel Core i7 / AMD Ryzen 7 (or equivalent) \\
\hline
Memory (RAM) & 16 GB DDR4 \\
\hline
Storage & 512 GB NVMe SSD \\
\hline
Operating System & Windows 11 / Linux (Ubuntu) \\
\hline
Web Browsers & Google Chrome (V8 engine) and Mozilla Firefox for cross-browser compatibility testing \\
\hline
Mobile Emulation & Google Chrome DevTools device mode for responsive design testing \\
\hline
Frontend Server & Vercel Edge Network (Serverless) \\
\hline
Database Server & Neon Cloud (Serverless PostgreSQL, dynamically scaled) \\
\hline
\end{tabular}
\end{table}

\subsection{Software Tools and Technologies}

Table~\ref{tab:software} presents the complete software stack used throughout the project, covering both client-side and server-side technologies.

\begin{table}[H]
\centering
\caption{Software tools and versions used in Stag.io}
\label{tab:software}
\begin{tabular}{|p{3.5cm}|p{2.5cm}|p{7.5cm}|}
\hline
\textbf{Tool / Technology} & \textbf{Version} & \textbf{Purpose} \\
\hline
Node.js & v20 LTS & Server-side JavaScript runtime \\
\hline
Next.js & 16.1.6 & React framework for frontend (App Router) \\
\hline
React & 19.2.3 & UI component library \\
\hline
Express.js & 5.2.1 & Backend REST API framework \\
\hline
TypeScript & 5.x & Static typing for frontend and backend \\
\hline
Tailwind CSS & v4 & Utility-first CSS framework \\
\hline
Framer Motion & v12 & Animation library for React \\
\hline
PostgreSQL & v16 (via Neon) & Relational database management system \\
\hline
Drizzle ORM & 0.45.1 & TypeScript-first ORM with migrations \\
\hline
Socket.IO & Latest & Real-time WebSocket communication \\
\hline
JSON Web Token & Latest & JWT-based authentication (access + refresh) \\
\hline
Vitest & Latest & Unit testing framework \\
\hline
Cloudinary & SDK v2 & Cloud-based image and file storage \\
\hline
Google Gemini & API v1 & AI-powered chatbot service \\
\hline
Cloudflare Turnstile & Latest & CAPTCHA bot protection \\
\hline
Visual Studio Code & Latest & Primary development IDE \\
\hline
Figma / Draw.io & --- & UI design and UML diagramming \\
\hline
\end{tabular}
\end{table}

\section{Key Implementation Details}

In this section, we present the implementation of the platform's most distinctive features: the SmartMatch recommendation engine, the real-time notification system, and the security mechanisms.

\subsection{SmartMatch Recommendation Engine}

The SmartMatch engine is the core differentiator of Stag.io. Unlike traditional platforms that display offers in chronological order, SmartMatch computes a relevance score for each offer based on four dimensions of the student's profile. The final composite score is computed using the following weighted formula:

\[
\text{Score}_{\text{final}} = 0.50 \times S_{\text{skills}} + 0.25 \times S_{\text{dept}} + 0.15 \times S_{\text{loc}} + 0.10 \times S_{\text{title}}
\]

where each component score is normalized to the range $[0, 1]$, and the final result is expressed as a percentage from 0 to 100.

\subsubsection{Skill Matching (50\% Weight)}

The skill matching function compares the student's declared skills against the offer's requirements and description. It supports both exact matching (when the full skill name appears in the text) and partial matching (when individual words of a multi-word skill appear). Partial matches are weighted at 70\% of the full match value, ensuring that broad skills like ``Machine Learning'' can still match offers that mention individual keywords.

\begin{figure}[H]
\begin{lstlisting}[language=JavaScript]
export function scoreSkills(
  studentSkills: string[],
  requirements: string,
  description: string,
): { score: number; matchedSkills: string[] } {
  if (studentSkills.length === 0) return { score: 0, matchedSkills: [] };
  const offerText = norm(requirements + ' ' + description);
  const matchedSkills: string[] = [];
  let totalWeight = 0;

  for (const skill of studentSkills) {
    const skillNorm = norm(skill);
    if (offerText.includes(skillNorm)) {
      matchedSkills.push(skill);
      totalWeight += 1.0;
    } else {
      const words = skillNorm.split(' ');
      const partialHits = words.filter(
        w => w.length > 2 && offerText.includes(w));
      if (partialHits.length > 0) {
        const ratio = partialHits.length / words.length;
        if (ratio >= 0.5) {
          matchedSkills.push(skill);
          totalWeight += ratio * 0.7;
        }
      }
    }
  }

  const reqTokens = requirements.split(/[,;/\n]+/)
    .map(t => t.trim()).filter(t => t.length > 1);
  const normFactor = Math.min(
    Math.max(reqTokens.length, 1), studentSkills.length);
  const score = matchedSkills.length > 0
    ? Math.min(totalWeight / normFactor, 1.0) : 0;
  return { score, matchedSkills };
}
\end{lstlisting}
\caption{Implementation of the skill matching function (50\% weight)}
\label{fig:code-skills}
\end{figure}

\subsubsection{Department Relevance (25\% Weight)}

The department scoring function maps academic departments to industry-relevant keywords. A pre-built dictionary covers 17 academic domains (Computer Science, Business, Medicine, Law, Design, etc.), each associated with 10--20 industry keywords. The function measures how many of these keywords appear in the offer text to determine the relevance between the student's field of study and the internship content. When no department is provided, the function returns a neutral score of 0.3 to avoid penalizing students with incomplete profiles.

\subsubsection{Location Matching (15\% Weight)}

The location scorer extracts city-level information from the student's university domain (e.g., extracting ``constantine'' from \texttt{univ-constantine2.dz}) and compares it against the offer's location and the company's headquarters. An exact city match returns a score of 1.0, remote offers receive 0.7 (as they are location-agnostic), and completely unmatched locations receive 0.1.

\subsubsection{Title Relevance (10\% Weight)}

The title relevance function checks whether the student's skills appear in the offer title, providing a heuristic for job-title alignment. The score is capped at 1.0 and normalized by the smaller of the student's skill count and three, ensuring that students with many skills are not unfairly disadvantaged.

\subsubsection{Orchestration and Ranking}

The main matching function orchestrates the entire pipeline: it retrieves the student profile and all active offers from the database, computes the four component scores for each offer, applies the weighted formula, sorts the results by descending score, and returns the top $n$ matches. The entire computation executes in under 100\,ms per request, meeting the performance requirements specified in Chapter~2.

\subsection{Real-Time Notification System}

The platform uses Socket.IO to deliver instant notifications to connected users. When a triggering event occurs (e.g., application status change, new application received, agreement pending validation), the backend persists the notification in the database and simultaneously pushes it to the user's browser through a dedicated WebSocket room.

\begin{figure}[H]
\begin{lstlisting}[language=JavaScript]
export async function sendNotification(
  userId: string,
  type: 'application_status_changed' | 'new_application'
    | 'agreement_needs_validation'
    | 'company_pending_approval'
    | 'university_pending_approval',
  title: string,
  message: string,
  relatedId?: string,
): Promise<NotificationData> {
  const [notification] = await db.insert(notifications).values({
    userId, type, title, message,
    relatedId: relatedId ?? null,
  }).returning();

  // Push real-time notification via WebSocket
  try {
    getIO().to(userId).emit('notification', {
      id: notification.id,
      userId: notification.userId,
      type: notification.type,
      title: notification.title,
      message: notification.message,
      isRead: notification.isRead,
      relatedId: notification.relatedId,
      createdAt: notification.createdAt.toISOString(),
    });
  } catch { /* Socket.io may not be initialized in tests */ }

  return data;
}
\end{lstlisting}
\caption{Real-time notification dispatch via Socket.IO}
\label{fig:code-notification}
\end{figure}

On the client side, the Socket.IO connection is authenticated using the JWT access token passed during the WebSocket handshake. Each user is automatically joined to a private room identified by their user ID, ensuring that notifications are delivered exclusively to the intended recipient.

\begin{figure}[H]
\begin{lstlisting}[language=JavaScript]
export function initSocket(httpServer: HttpServer): Server {
  io = new Server(httpServer, {
    cors: {
      origin: (process.env.CLIENT_URL || 'http://localhost:3000')
        .replace(/\/+$/, ''),
      credentials: true,
    },
  });

  io.on('connection', (socket) => {
    const token = socket.handshake.auth.token;
    if (!token) { socket.disconnect(); return; }

    try {
      const payload = jwt.verify(
        token, process.env.JWT_SECRET!) as JwtPayload;
      socket.join(payload.sub); // Join user-specific room
    } catch {
      socket.disconnect();
    }
  });

  return io;
}
\end{lstlisting}
\caption{WebSocket server initialization with JWT-authenticated rooms}
\label{fig:code-socket}
\end{figure}

\subsection{Security Implementation}

The security layer of Stag.io implements several defense-in-depth mechanisms to protect user accounts and API endpoints.

\subsubsection{Account Lockout Mechanism}

To prevent brute-force attacks, the system implements an in-memory account lockout tracker. After five consecutive failed login attempts, the account is locked for 15 minutes. The counter is automatically reset upon successful authentication. A periodic cleanup routine runs every five minutes to evict expired lockout records from memory.

\begin{figure}[H]
\begin{lstlisting}[language=JavaScript]
const MAX_ATTEMPTS = 5;
const LOCKOUT_DURATION_MS = 15 * 60 * 1000; // 15 minutes

const attempts = new Map<string, AttemptRecord>();

export function recordFailedAttempt(email: string): boolean {
  const key = email.toLowerCase();
  const record = attempts.get(key)
    ?? { count: 0, lockedUntil: null };
  record.count += 1;

  if (record.count >= MAX_ATTEMPTS) {
    record.lockedUntil = Date.now() + LOCKOUT_DURATION_MS;
    attempts.set(key, record);
    return true; // Account is now locked
  }

  attempts.set(key, record);
  return false;
}

export function isAccountLocked(email: string): number {
  const key = email.toLowerCase();
  const record = attempts.get(key);
  if (!record?.lockedUntil) return 0;

  const remaining = record.lockedUntil - Date.now();
  if (remaining <= 0) { attempts.delete(key); return 0; }
  return Math.ceil(remaining / 1000);
}
\end{lstlisting}
\caption{Account lockout mechanism for brute-force protection}
\label{fig:code-lockout}
\end{figure}

\subsubsection{JWT Token Management and Blacklisting}

The authentication flow relies on short-lived access tokens (15 minutes) and long-lived refresh tokens (7 days) with full token rotation on each refresh. Upon logout, the current access token is added to an in-memory blacklist indexed by its expiration time, preventing its reuse during the remaining validity window. A periodic cleanup routine removes expired entries to prevent memory leaks.

\section{Database Structure and Data Implementation}

The application handles a multitude of interconnected data sets governing users, internships, and dynamic features like notifications and reviews.

\subsection{Database Structure}

The database structure maps exactly to the relational model envisioned during the design phase in Chapter~3. It utilizes PostgreSQL as the primary persistence engine, hosted on Neon Cloud for serverless scalability. The schema heavily depends on relational constraints (\texttt{FOREIGN KEY}, \texttt{UNIQUE}, \texttt{ON DELETE CASCADE}) to naturally maintain data integrity between companies, their offers, and the student applications.

\subsection{Implementation of Data and Processing}

To manipulate this data structure from our backend, we implemented Drizzle ORM --- a TypeScript-first ORM that provides compile-time type safety across all database queries. Below is an example snippet showing how the \texttt{users} and \texttt{internship\_offers} data structures are translated into pure TypeScript schema definitions.

\begin{figure}[H]
\begin{lstlisting}[language=JavaScript]
import { pgTable, text, boolean, timestamp, pgEnum }
  from 'drizzle-orm/pg-core';
import { createId } from '@paralleldrive/cuid2';

export const userRoleEnum = pgEnum('UserRole',
  ['student', 'company', 'admin', 'university']);

export const users = pgTable('users', {
  id: text('id').primaryKey()
    .$defaultFn(() => createId()),
  email: text('email').notNull().unique(),
  passwordHash: text('passwordHash').notNull(),
  role: userRoleEnum('role').default('student').notNull(),
  isEmailVerified: boolean('isEmailVerified')
    .default(false).notNull(),
  createdAt: timestamp('createdAt', { withTimezone: true })
    .defaultNow().notNull(),
});

export const internshipOffers = pgTable('internship_offers', {
  id: text('id').primaryKey()
    .$defaultFn(() => createId()),
  companyId: text('companyId').notNull()
    .references(() => companies.id, { onDelete: 'cascade' }),
  title: text('title').notNull(),
  status: text('status').default('active').notNull(),
  createdAt: timestamp('createdAt', { withTimezone: true })
    .defaultNow().notNull(),
});
\end{lstlisting}
\caption{Data modeling with Drizzle ORM (TypeScript schema definition)}
\label{fig:code-schema}
\end{figure}

This schema-as-code approach ensures that any change to the database structure is versioned, reviewable, and automatically reflected in the TypeScript types used throughout the application.

\section{Testing and Validation}

To guarantee the correctness and reliability of the platform's critical components, we adopted a structured unit testing strategy using the Vitest testing framework. This section presents the tests implemented for the SmartMatch engine and the security modules.

\subsection{SmartMatch Engine Tests}

The SmartMatch engine was validated through a comprehensive test suite covering each of the four scoring dimensions individually, as well as the weighted final score formula. Table~\ref{tab:test-smart} summarizes the key test scenarios.

\begin{table}[H]
\centering
\caption{Summary of SmartMatch unit test scenarios}
\label{tab:test-smart}
\footnotesize
\begin{tabular}{|p{3cm}|p{5.5cm}|p{4.5cm}|}
\hline
\textbf{Function} & \textbf{Test Scenario} & \textbf{Expected Result} \\
\hline
\texttt{scoreSkills} & Empty student skills & Score = 0, no matches \\
\hline
\texttt{scoreSkills} & Exact skill names in requirements & Skills matched, score $>$ 0 \\
\hline
\texttt{scoreSkills} & Skills found only in description & Skills matched, score $>$ 0 \\
\hline
\texttt{scoreSkills} & Completely unrelated skills & Score = 0 \\
\hline
\texttt{scoreSkills} & Multi-word partial match & Skill matched via partial \\
\hline
\texttt{scoreSkills} & Case-insensitive input & Matches regardless of case \\
\hline
\texttt{scoreDept} & Null department & Neutral score (0.3) \\
\hline
\texttt{scoreDept} & CS student + software offer & Score $\geq$ 0.5 \\
\hline
\texttt{scoreDept} & Medicine + web dev offer & Score $<$ 0.3 \\
\hline
\texttt{scoreLoc} & University city matches offer & Score = 1.0 \\
\hline
\texttt{scoreLoc} & Remote offer & Score = 0.7 \\
\hline
\texttt{scoreLoc} & No location match & Score = 0.1 \\
\hline
\texttt{scoreTitle} & Skills appear in title & Score $>$ 0.5 \\
\hline
Weighted & Weights sum to 1.0 & $0.50+0.25+0.15+0.10$ \\
\hline
Weighted & Full scenario (CS / Constantine) & Score $\in [50, 100]$ \\
\hline
\end{tabular}
\end{table}

The following excerpt illustrates the test implementation for the skill matching function:

\begin{figure}[H]
\begin{lstlisting}[language=JavaScript]
describe('scoreSkills()', () => {
  it('returns 0 for empty student skills', () => {
    const result = scoreSkills(
      [], 'React, Node.js', 'Build frontend apps');
    expect(result.score).toBe(0);
    expect(result.matchedSkills).toEqual([]);
  });

  it('matches exact skill names in requirements', () => {
    const result = scoreSkills(
      ['React', 'Node.js'],
      'Looking for React and Node.js developer',
      'Full stack web application',
    );
    expect(result.matchedSkills).toContain('React');
    expect(result.matchedSkills).toContain('Node.js');
    expect(result.score).toBeGreaterThan(0);
  });

  it('is case-insensitive', () => {
    const result = scoreSkills(
      ['REACT', 'node.JS'],
      'We use react and node.js', '',
    );
    expect(result.matchedSkills).toContain('REACT');
    expect(result.matchedSkills).toContain('node.JS');
  });
});
\end{lstlisting}
\caption{Unit tests for the SmartMatch skill matching function}
\label{fig:code-test-skills}
\end{figure}

\subsection{Security Module Tests}

The token blacklist and login limiter modules were also validated through unit tests to ensure correct behavior under normal and edge-case conditions.

\begin{figure}[H]
\begin{lstlisting}[language=JavaScript]
describe('Token Blacklist', () => {
  it('a fresh token is not blacklisted', () => {
    expect(isTokenBlacklisted('some-random-token-abc123'))
      .toBe(false);
  });

  it('a blacklisted token is detected', () => {
    const token = 'test-token-' + Date.now();
    const futureExpiry =
      Math.floor(Date.now() / 1000) + 3600;
    blacklistToken(token, futureExpiry);
    expect(isTokenBlacklisted(token)).toBe(true);
  });

  it('different tokens are independent', () => {
    const token1 = 'token-one-' + Date.now();
    const token2 = 'token-two-' + Date.now();
    const futureExpiry =
      Math.floor(Date.now() / 1000) + 3600;
    blacklistToken(token1, futureExpiry);
    expect(isTokenBlacklisted(token1)).toBe(true);
    expect(isTokenBlacklisted(token2)).toBe(false);
  });
});
\end{lstlisting}
\caption{Unit tests for the JWT token blacklist module}
\label{fig:code-test-blacklist}
\end{figure}

All tests execute in under 200\,ms and can be run via the command \texttt{npx vitest run}. A total of 20 test cases cover the SmartMatch engine and 6 test cases cover the security modules, achieving full branch coverage for these critical components.

\section{User Interfaces and Software Manipulation}

In this section, we present the practical use and manipulation of the software. The user interface was built using React 19 with Tailwind CSS v4 and Framer Motion v12 to provide maximum clarity, smooth animations, and a premium user experience. The interface adapts intelligently to the actor's platform role.

\subsection{Public Portal and Authentication}

The initial touchpoint for all users. The public portal presents the internship catalog, company listings, and platform information pages. It guides unauthenticated visitors toward secure role-based registration and login.

\begin{figure}[H]
\centering

\includegraphics[width=\textwidth,keepaspectratio]{ui-homepage.png}
\caption{Homepage of the Stag.io platform, allowing open discovery}
\label{fig:ui-home}
\end{figure}

\begin{figure}[H]
\centering
\fbox{\rule{0pt}{3in} \rule{0.8\textwidth}{0pt}} % Placeholder box
\caption{Authentication interface equipped with CAPTCHA protection for safe manipulation}
\label{fig:ui-login}
\end{figure}

\subsection{Student Dashboard and SmartMatch}

Once logged in, students access a personalized dashboard displaying application statistics, saved offers, and recent activity. From this dashboard, they can activate the SmartMatch engine to receive offers ranked by relevance to their profile.

\begin{figure}[H]
\centering
\fbox{\rule{0pt}{3in} \rule{0.8\textwidth}{0pt}} % Placeholder box
\caption{Student dashboard showing application statistics and saved offers}
\label{fig:ui-student-dashboard}
\end{figure}

\begin{figure}[H]
\centering
\fbox{\rule{0pt}{3in} \rule{0.8\textwidth}{0pt}} % Placeholder box
\caption{SmartMatch recommendation engine generating internship matches with relevance scores}
\label{fig:ui-smartmatch}
\end{figure}

\subsection{Company and University Dashboards}

The company dashboard provides tools for managing internship publications, reviewing incoming applications, and tracking recruitment statistics. The university dashboard enables institutional oversight of affiliated students and internship agreement validation.

\begin{figure}[H]
\centering
\fbox{\rule{0pt}{3in} \rule{0.8\textwidth}{0pt}} % Placeholder box
\caption{Company interface for managing internship publications and reviewing application pipelines}
\label{fig:ui-company-dashboard}
\end{figure}

\begin{figure}[H]
\centering
\fbox{\rule{0pt}{3in} \rule{0.8\textwidth}{0pt}} % Placeholder box
\caption{University portal for tracking student milestones and validating internship agreements}
\label{fig:ui-university-dashboard}
\end{figure}

\section{Conclusion}

This chapter encapsulated the practical realization of the Stag.io platform. We began by outlining the proposed system architecture through component and deployment diagrams, then identified the hardware and software resources required for development and deployment.

The core of this chapter focused on the implementation of the platform's key features: the SmartMatch recommendation engine with its four-dimensional scoring algorithm, the real-time notification system powered by Socket.IO, and the security mechanisms including account lockout and JWT token blacklisting. Each feature was documented with actual source code excerpts from the project repository.

We then presented the testing strategy, demonstrating how unit tests validate the correctness of the SmartMatch engine and security modules using the Vitest framework. Finally, we showcased the user interfaces that govern day-to-day platform interaction across all actor roles.
